% RUN / ! -- os/samir/cmd/mutate.c d_run (Using III+ U5-185/186, PDF
% 297-298: RUN/! <command>; the rest of the line goes verbatim to the
% command processor, which performs any redirection, then control returns;
% one CR LF before and one after the child; 92 when the processor cannot be
% loaded).  The OS PAL (os/samir/pal/pal_milton.c) EXECs COMSPEC with
% " /C <line>"; the example's command is lower-case ECHO with a redirection,
% which InitechDOS COMMAND.COM and the host processor of the replay both run
% and whose file TYPE shows the same (the output to the screen would not be
% part of the replayed transcript).  Lane samirfix, bead initech-n8q8.
% Gates: test-samir-dml (run), test-samir-dml-emu.
\begin{XbaseEntry}{RUN}{command}\XAlso{!}{!}%
Carries out an InitechDOS command or runs a program, then returns to
InitechBase.
\Syntax
\SyntaxLine{\kw{RUN} \xp{command}}
\par or\par
\SyntaxLine{\kw{!} \xp{command}}
\Usage
\cmd{RUN} passes the rest of the line, as it stands, to the InitechDOS
command processor, \file{COMMAND.COM}, which carries it out: an internal
command such as \texttt{DIR}, \texttt{COPY} or \texttt{ECHO}, a program, or a
batch file.  Redirection with \texttt{>} is carried out by the command
processor.  When the command ends, InitechBase goes on at the dot prompt or
in the program.  Open databases stay open, and the record pointer is where it
was.  \cmd{!} is another name for \cmd{RUN}.

\file{COMMAND.COM} is found through the COMSPEC setting of the environment,
or else on drive A.  If it cannot be loaded, \cmd{RUN} reports
\xmsg{92}{Unable to load COMMAND.COM.}  The command processor and the
program must fit in the memory InitechBase leaves free.
\Tips
Do not \cmd{RUN} a program that stays resident in memory: the memory it keeps
is not returned to InitechBase.
\Examples
A note is written to a text file by the InitechDOS \texttt{ECHO} command and
displayed with \cmd{TYPE}.  The record pointer is unchanged:
\DotExample{examples/c-run}
\SeeAlso{TYPE, QUIT}
\end{XbaseEntry}
